\chapter{Introduction}\label{ch:introduction}

\section{Motivation}\label{sec:intro:motivation}

A robot that plans with a generative model buys expressiveness with computation. Diffusion
planners \parencite{janner2022planning,chi2023diffusion} represent multi-modal demonstration data
\parencite{jia2024towards} well and, when paired with a constrained projection step
\parencite{romer2025diffusion}, can respect dynamics and safety at deployment. The price is the sampling
loop: producing one plan costs tens to hundreds of network evaluations --- the \ac{NFE} of the
sampler --- and in a receding-horizon controller that cost is paid again at every timestep.

This cost is not incidental but \emph{structural}, and the structure is what makes it
addressable. In a diffusion model the number of steps is a
property of the trained noise schedule \parencite{ho2020denoising}; reducing it after training
degrades the sample, and recovering quality needs a different sampler
\parencite{song2021denoising}. In a deterministic transport model \parencite{lipman2023flow} the same
quantity is a property of the \ac{ODE} solver, and is chosen at inference. This thesis moves the
step budget from training time to inference time and measures what the move costs.

Once sampling is cheap, constraint enforcement becomes the dominant cost, and how constraints are
enforced becomes a question in its own right.

\section{Problem Statement}\label{sec:intro:problem}

The object of study is a closed-loop controller in which, at each timestep, a conditional
generative model proposes a trajectory over a fixed horizon, a constraint mechanism makes
that trajectory admissible, and the first action is executed. This thesis holds that
architecture fixed and varies two parts of it: the generative model and the projection method.

\emph{Flow matching} in the title names the family of generative models that transport noise to data
along an ordinary differential equation. The three generative models studied here belong to it and
differ in the velocity their network learns: \emph{instantaneous-velocity matching},
\emph{analytic average-velocity matching} and \emph{consistency-interpolated average-velocity
matching}. The study covers two observation modalities, state and camera images, and two embodiments, a manipulator
and a quadrotor.

The generative models and projection methods are first compared with the diffusion baseline on the
obstacle-avoidance benchmark of \ac{DPCC} \parencite{romer2025diffusion,jia2024towards}, with its
setup unchanged, and are then tested in a vision-conditioned alignment task and a quadrotor benchmark
built for this thesis. \autoref{ch:setup} describes the evaluation.

\section{Research Questions}\label{sec:intro:rqs}

\begin{description}
  \item[RQ1 --- Generative model.] Under an identical backbone, projection, dataset and evaluation
    code, does flow matching --- a learned velocity field, the transport \ac{ODE} it generates, the
    step budget that becomes a solver argument, and the average-velocity objectives built on that
    field --- improve the quality--cost frontier of constrained predictive control over the diffusion
    model it replaces, and does the improvement order the three objectives?
  \item[RQ2 --- Constraints.] Given one fixed feasible set and one solver, at a step budget large
    enough for projection to guide the sampler, does endpoint projection reach the goal with fewer
    constraint violations, fewer control steps or less computation than the per-step (iterate)
    projection of \ac{DPCC}, and below which budget does the comparison cease to be well-posed?
  \item[RQ3 --- Transfer.] Obstacle avoidance is planar, observed as a state, kinematic at the
    planning level, and solvable by a classical planner with a collision constraint. Do the answers
    to RQ1 and RQ2 hold when the observation becomes camera images, when the task requires aligning
    an object through contact, and when the platform is an underactuated, open-loop unstable
    quadrotor moving through three-dimensional obstacle fields --- settings whose dynamics and
    perception are closer to those of a physical deployment?
\end{description}

\srcnote{TARGET \S7 maps each RQ to its evidence sections; the former RQ2 on the step budget is
folded into RQ1 (v2.22).}

\section{Contributions}\label{sec:intro:contributions}

This thesis makes six contributions.

\begin{enumerate}
  \item \textbf{Flow matching in place of diffusion.} The diffusion model of \ac{DPCC}
    \parencite{romer2025diffusion} is replaced by instantaneous-velocity matching, the flow-matching
    objective \parencite{lipman2023flow,lipman2024guide}, which learns a velocity field and generates a
    trajectory by integrating an ordinary differential equation. Backbone, projection, data and
    evaluation protocol are kept, and the number of network evaluations becomes a choice made at
    deployment. The comparison is made first on
    the obstacle-avoidance benchmark of \ac{DPCC}, with its setup unchanged.
    \hole{Result sentence for this contribution --- from Ch.~6, supplied by v3 as a \texttt{For v2:} line.}
  \item \textbf{Analytic average-velocity matching.} Analytic average-velocity matching
    \parencite{geng2025meanflow} learns the average velocity over an interval instead of the
    instantaneous velocity, with a regression target constructed analytically through a
    Jacobian--vector product, so that a single network evaluation can cover the whole integration. It
    replaces instantaneous-velocity matching within the same architecture.
    \hole{Result sentence for this contribution --- from Ch.~6, supplied by v3 as a \texttt{For v2:} line.}
  \item \textbf{Consistency-interpolated average-velocity matching.} Consistency-interpolated
    average-velocity matching \parencite{zhang2025alphaflow} also learns the average velocity, with a
    target interpolated by a consistency step ratio between the instantaneous velocity and a
    stop-gradient prediction of the network itself; the ratio is annealed during training. Both
    average-velocity objectives achieve state-of-the-art one- and two-step image generation
    \parencite{geng2025meanflow,zhang2025alphaflow}; this thesis evaluates both as planners inside a
    constrained controller.
    \hole{Result sentence for this contribution --- from Ch.~6, supplied by v3 as a \texttt{For v2:} line.}
  \item \textbf{Endpoint projection against DPCC's per-step projection.} The central contribution of
    DPCC is a model-based projection applied to each denoising step \parencite{romer2025diffusion}.
    Endpoint projection \parencite{li2025hardflow} instead projects the trajectory the sampler is
    predicted to reach and steers the sampler towards it. Both are implemented on the same feasible set
    and solver and compared, together with the step budget below which endpoint projection reduces to
    projection after sampling. On obstacle avoidance the guiding step count is zero at one and two
    network evaluations, where endpoint projection reduces to projecting the finished sample; above that
    budget, neither projection method leads consistently (\autoref{sec:res:constraints:degenerate}).
    Where it has guiding steps, endpoint projection takes less time per control step than per-step
    projection in every environment; on UAV-corridor it is also less often collision-free
    (\autoref{sec:res:constraints}).
    \srcnote{Result sentences from v3 Ch.~6 (\texttt{cross\_draft/to\_v2/FROM\_v3\_20260915\_result\_sentences.md},
    worded after \texttt{sec:res:avoiding:projection} and \texttt{tab:hf-ladder} as of v3.16; cost sentence
    from \texttt{FROM\_v3\_20260917\_projection\_cost\_resolved.md}, v3.18).}
  \item \textbf{A vision-conditioned alignment task.} Once established on the state-based benchmark
    of \ac{DPCC}, all generative models and both projection methods are tested on a more complex
    task built for this thesis from the alignment task of D3IL, in which the robot observes two camera
    views. The generative network is given a visual encoder, and the projection is given the dynamics
    rows and the constraint set of the new task.
    \hole{Result sentence for this contribution --- from Ch.~6, supplied by v3 as a \texttt{For v2:} line.}
  \item \textbf{A quadrotor benchmark.} The controller is carried to an underactuated, open-loop
    unstable system: the manipulator with inverse kinematics is replaced by a quadrotor commanded
    through its rotor thrusts, and the stack is tested from state observations in three environments
    of increasing difficulty. Apart from the vehicle model, everything in this benchmark was built for
    this thesis.
    \begin{enumerate}
      \item The quadrotor model is the Skydio X2 from MuJoCo Menagerie
        \parencite{todorov2012mujoco,zakka2022menagerie}.
      \item It is controlled by a cascaded geometric tracking controller \parencite{lee2010geometric},
        with sampling-based MuJoCo predictive control \parencite{howell2022predictive} as a baseline
        controller.
      \item The three environments, their constraint sets and the expert demonstrations are built
        for this thesis in MuJoCo \parencite{todorov2012mujoco}, inside the environment code of
        D3IL.
    \end{enumerate}
    \hole{Result sentence for this contribution --- from Ch.~6, supplied by v3 as a \texttt{For v2:} line.}
\end{enumerate}

\section{Outline}\label{sec:intro:outline}

\hole{Write last. One paragraph, one sentence per chapter.}

% =============================================================================
